% ============================================================================
% Camera-ready replacement text for ICICF-2027 Paper 768 (IEEEtran).
% Drop each block into the matching section of the manuscript.
% Every number below is produced by paper_revision/code/run_all.sh.
% Lines marked %CONFIRM must be checked by the authors.
% Figures: paper_revision/figures/*.png
% ============================================================================

% ---------------------------------------------------------------- ABSTRACT
\begin{abstract}
Laboratory motion capture for sprint biomechanics costs USD 100k--500k, which
puts it out of reach for most athletes. We present an end-to-end, reproducible
MLOps pipeline that combines a single 60\,fps smartphone video (MediaPipe Pose)
with a consumer smartwatch IMU (100\,Hz accelerometer and gyroscope) to extract
19 sprint biomechanics features. Performance labels come from an objective,
automatically measured quantity: each athlete's maximal running velocity,
normalised by body size, rather than subjective ratings. The integration of the
two sensors is validated quantitatively. Stride frequency from video and from
the watch agree with a mean absolute error of 0.07\,Hz ($r=0.95$), and
clock-consistency checks automatically flagged two of nine watch recordings
that did not belong to the filmed sprint. On commodity hardware the pipeline
processes a 2--4\,s sprint clip in 2.1\,s, and the exported 52\,KB TFLite
classifier runs in microseconds. In a leakage-free leave-one-subject-out pilot
study with 8 athletes (9 trials), three-tier performance classification did
not exceed chance with statistical significance (best 55.6\%, permutation
$p=0.16$). Accuracies previously reported for this system arose from data
leakage, and we document how it occurred. The pipeline is versioned (DVC, MLflow), containerised and covered by 43 automated tests, with a quality
gate that quarantines mismatched sensor files, a promotion gate that refuses unvalidated models, and drift
monitoring whose power analysis shows that at least 100 reference recordings are needed before alarms can be armed. The pipeline, measurements and
negative result define a validated starting point and a concrete
data-collection target for low-cost sprint analytics.
\end{abstract}

% ------------------------------------------------- CONTRIBUTIONS (Sec. I end)
The contributions of this work are:
(1)~a reproducible smartphone + smartwatch pipeline for sprint biomechanics
with experiment tracking (MLflow), configuration-driven training and a deployed
web interface;
(2)~an objective, annotation-free performance label (body-size-normalised
maximal velocity) that removes subjective rating from the loop and enables
automatic labelling of new data;
(3)~a quantitative validation of video--IMU integration, including automatic
detection of mismatched sensor files;
(4)~measured, per-stage latency on two hardware configurations;
(5)~a leakage-free pilot evaluation, with a power analysis that turns the
negative result into a concrete data requirement for drift-triggered
retraining.

% ---------------------------------------------- III-B DATA COLLECTION (replace)
\subsection{Sprint Dataset}
Nine maximal-effort fly sprints were recorded from \textbf{eight} athletes (one
athlete ran two trials) in a single outdoor session on a synthetic track on
25~March~2026. %CONFIRM: ages, sex, recruitment, ethics approval/consent wording
Video was captured with a tripod-mounted Samsung Android smartphone %CONFIRM model
at 1920$\times$1080, 60\,fps (HEVC), with the camera orthogonal to the running
direction so that athletes crossed the field of view in 0.9--1.7\,s.
A Samsung Galaxy Watch %CONFIRM model and wrist
recorded triaxial acceleration and angular velocity at a measured
99.3--99.6\,Hz using the Sensor Logger application. Seven recordings share one
camera position and two share a second position; the velocity normalisation in
Section~\ref{sec:labels} makes measurements comparable across both.

% ------------------------------------------------- III-C LABELS (new section)
\subsection{Objective Performance Labels}\label{sec:labels}
Each recording is labelled from a quantity that is \emph{not} a model input:
the athlete's horizontal velocity. The mid-hip landmark's horizontal position
is fitted against time with a Theil--Sen estimator ($R^2\geq0.998$ for all
recordings). The pixel velocity is divided by a posture-invariant body length
$L$ (trunk + thigh + shank, median over frames), which corresponds to
$0.779\times$ stature according to standard anthropometric ratios. This yields
velocity in statures per second, independent of camera distance. Assuming a
stature of 1.70\,m, velocities range from 6.7 to 9.4\,m/s. Recordings are
assigned to cohort tertiles (3 per tier): Beginner $<7.45$\,m/s $\leq$
Intermediate $<7.79$\,m/s $\leq$ Elite. Both trials of the repeated athlete
fall in the same tier. Because the label is measured automatically, new
recordings can be labelled without expert annotation. MediaPipe occasionally
locked onto a static false detection before the runner entered the frame, so
only the continuous track with the largest horizontal travel is retained.

% ---------------------------------------------- III-D FEATURES (replace lists)
\subsection{Feature Extraction}
\textbf{Video (9 features).} Mean knee and hip angles (2-D, Eq.~(1)); knee and
hip symmetry indices (Eq.~(2)); trunk lean; vertical hip oscillation
(detrended, in statures); step frequency, ground contact time (GCT) and duty
factor. Stance is detected when the slower ankle's horizontal speed falls below
$0.4\,v$ (where $v$ is the runner's speed), because a planted foot is
stationary in the world frame. This threshold was set so that all recordings
yield physiologically plausible step rates (3.2--4.6\,Hz) and GCT
(70--117\,ms); it was not tuned on classification accuracy. Stride length,
clip duration and frame count are excluded: the first is algebraically
velocity divided by step rate (label leakage), and the others depend on how
each clip was trimmed.
\textbf{IMU (10 features).} Peak, mean and SD of acceleration and angular
velocity magnitude; RMS jerk (Eq.~(4)); arm-swing frequency (Welch PSD of the
gyroscope); and SPARC smoothness of acceleration and angular velocity.

% ------------------------------------------- III-E INTEGRATION (replace fusion)
\subsection{Video--IMU Integration}
The watch and phone clocks are independent. The IMU segment used is the
highest-intensity window whose length equals the time the runner is tracked in
video. Accelerometer and gyroscope streams are merged by nearest timestamp
(20\,ms tolerance). Fusion is early: the 9 video and 10 IMU features are
concatenated after per-fold standardisation. Integration is validated in three
ways (Fig.~\ref{fig:sync}).
(i)~\emph{Clock consistency:} for 7 of 9 recordings, the video creation time
follows the watch's last sample by 1.6--5.7\,s (median 2.1\,s). Two recordings
differ by $-102$\,s and $-422$\,s, meaning those watch files were captured
during a different run; they are flagged, and Table~\ref{tab:results} includes
a sensitivity analysis without them.
(ii)~\emph{Cadence agreement} on the 7 consistent recordings: stride frequency
from the video wrist landmark versus the watch gyroscope gives MAE 0.073\,Hz,
bias $-0.05$\,Hz and $r=0.95$.
(iii)~\emph{Waveform agreement:} the normalised cross-correlation between video
wrist velocity and the dominant gyroscope axis has $|r|$ of 0.51--0.94
(median 0.76).

% ------------------------------------------- III-F PROTOCOL (replace splitting)
\subsection{Evaluation Protocol}
All models are evaluated with leave-one-subject-out (LOSO) cross-validation (8
folds; both trials of the repeated athlete are always held out together).
Within each fold, and using the training athletes only, a StandardScaler is
fitted and the training set is optionally augmented to 300 samples per class:
Gaussian jitter at 5\% of feature SD, plus same-class linear interpolation
($\alpha\sim U(0.1,0.9)$). The held-out athlete never contributes to scaling,
augmentation or early stopping. Models: logistic regression ($C=1$), RBF-SVM
($C=1$), random forest (100 trees, depth 5), and the MLP of Fig.~4
(128--64--32, ReLU, batch normalisation, dropout 0.4/0.3/0.2, L2 $10^{-3}$,
Adam $10^{-3}$, batch 32, at most 100 epochs, early stopping with patience 10
on a 15\% validation split of the training fold, 5 seeds). We report accuracy,
balanced accuracy, macro precision, recall and F1 over the 9 pooled
out-of-fold predictions; 95\% CIs from 2000 subject-level bootstrap resamples;
and label-permutation $p$-values (200 permutations, classical models only).

% ---------------------------------------------------- IV RESULTS (replace)
\begin{table}[t]
\caption{LOSO classification (chance = 33.3\%). Best model per modality.}
\label{tab:results}
\centering
\begin{tabular}{llcccc}
\hline
Modality & Model & Acc. & Macro F1 & 95\% CI & $p$ \\
\hline
\multicolumn{6}{l}{\emph{All recordings (8 athletes, 9 trials)}}\\
Video & MLP & 35.6$\pm$16.5 & 0.31 & 13--75 & -- \\
IMU   & LogReg+aug & 55.6 & 0.48 & 25--88 & 0.16 \\
Fused & MLP & 28.9$\pm$16.9 & 0.24 & 13--89 & -- \\
\multicolumn{6}{l}{\emph{Watch files verified (6 athletes, 7 trials)}}\\
Video & SVM+aug & 28.6 & 0.33 & 0--67 & 0.44 \\
IMU   & RF+aug & 57.1 & 0.52 & 25--86 & 0.10 \\
Fused & RF+aug & 14.3 & 0.10 & 0--33 & 0.70 \\
\hline
\end{tabular}
\end{table}

\subsection{Classification Results}
Table~\ref{tab:results} and Fig.~\ref{fig:cm} report the leakage-free results.
No configuration is significantly above chance (smallest permutation
$p=0.10$), and early fusion does not outperform the single modalities. The
paper's MLP with augmentation reaches 26.7\% (video), 51.1\% (IMU) and 0\%
(fused). Several configurations score below chance, a known artefact of
leave-one-out validation with balanced classes of size 3, where removing a
test sample biases training against its class. A LOSO ridge regression of
velocity also fails to beat a training-mean predictor (MAE 0.80--0.88\,m/s
vs.\ 0.71\,m/s). Descriptively, GCT correlates negatively ($r=-0.52$) and step
frequency positively ($r=0.40$) with velocity, consistent with sprint
literature~[15], but $n=9$ is too small for inference.

\emph{Relation to previously reported accuracy.} An earlier version of this
system reported 82--100\% accuracy. Those figures came from (a) labels derived
by K-Means clustering of the input features themselves, and (b) synthetic
samples generated from all athletes before the train/test split, so that every
test athlete had near-copies in the training data. Removing both effects gives
the results above. We report this explicitly because the same pitfalls are
common in small-sample sports analytics.

% --------------------------------------------- IV-E LATENCY (replace)
\begin{table}[t]
\caption{Measured latency, median (p95). MediaPipe 0.10.9, TF 2.18.1.}
\label{tab:latency}
\centering
\begin{tabular}{lcc}
\hline
Stage & M5, all cores & M5, E-cores only \\
\hline
Decode / frame & 1.4\,ms & 10.6\,ms \\
Pose lite / frame & 7.1 (13.0)\,ms & 45.8 (75.8)\,ms \\
Pose full / frame$^\dagger$ & 11.5 (17.9)\,ms & 61.3 (114.2)\,ms \\
Pose heavy / frame & 49.3 (56.3)\,ms & 212.3 (273.3)\,ms \\
Features / sprint & 8.1\,ms & 49.9\,ms \\
MLP, Keras \texttt{predict} & 14.5\,ms & 95.3\,ms \\
MLP, TFLite (52\,KB) & 1.1\,$\mu$s & 9.3\,$\mu$s \\
End-to-end / clip$^\dagger$ & 2.1\,s & 12.7\,s \\
\hline
\multicolumn{3}{l}{$^\dagger$deployed configuration (model complexity 1)}
\end{tabular}
\end{table}
Pose estimation dominates the runtime (Table~\ref{tab:latency}). Once the
classifier is exported to TFLite, its cost is negligible; the Keras
\texttt{predict} overhead (14.5--95\,ms) is framework overhead, not model
compute. On-phone execution was not measured and is left to future work.

% --------------------------------------------- V-D MLOPS SYSTEM (new, replaces drift section)
\subsection{MLOps System}
\emph{Architecture.} A DVC pipeline (pose $\rightarrow$ features $\rightarrow$ quality gate $\rightarrow$ training)
records the raw-data hash; every MLflow run logs the git commit, data hash, parameters, leave-one-subject-out (LOSO)
metrics and report. A model is served only if it is the registry's \texttt{champion}. A challenger is promoted only if
it is significantly above chance (permutation $p<0.05$, max-statistic over all candidates) and its macro-F1 is not below
the champion's; the first model is registered as \emph{provisional} and the API labels it as such. The FastAPI service
(\texttt{/predict}, \texttt{/health}, \texttt{/metrics}, \texttt{/drift}) is packaged in Docker; each request's features
go to a feature store used by the drift monitor and the retraining job. 43 automated tests cover velocity recovery,
leakage, the quality gate, drift false-alarm and power, the promotion gate, registry restore and the API contract.

\emph{Data-quality gate.} Seven automatic checks (frame rate, IMU rate, runner-track fit, tracked frames, finite features,
video--IMU cadence agreement, clock consistency) quarantine Chandra1 and Jeeva without manual inspection.

\emph{Drift thresholds.} Each feature is tested with a two-sample KS test at the Bonferroni-corrected
$\alpha=0.05/19=0.0026$ (family-wise false-alarm rate $\leq5\%$; 300 null windows gave $\leq10\%$ in testing); PSI is a
secondary signal. The power analysis (Fig.~\ref{fig:power}) shows that a 1-SD shift is detected with 84\% probability
only with a reference of $\geq100$ and a window of $\geq30$; alarms are therefore \emph{armed} only above those sizes
and are informational below them. Our reference currently has 7 recordings, so no alarm is armed.
Retraining is triggered by an armed alarm, by every 3 new quality-passed recordings, or manually.

\emph{Device-shift replay} (7 real clips through the API): an unmodified replay raises no alarm; a 50\,Hz watch is
rejected by the quality gate in 7/7 cases; a $\times1.5$ accelerometer gain raises an informational alarm on
\texttt{peak\_accel}, \texttt{std\_accel} and \texttt{rms\_jerk} ($p=0.00058$).

\emph{Retraining under a changing population.} Replaying the nine recordings in capture order from a cold start, two
recordings are quarantined, a first (provisional, not significant: $p=0.099$) model appears at the 8th arrival, and the
retrained challenger at the 9th is refused ($p=0.65$, macro-F1 0.24 vs.\ 0.56). The loop and gate therefore work, but
\textbf{an accuracy benefit from retraining could not be demonstrated} at this sample size.

\begin{table}[t]
\caption{Load and fault tests (Docker, 10-core node, real 2--3\,s clips including upload). Zero failures in 214 requests.}
\label{tab:ops}\centering\small
\begin{tabular}{ccccc}
\hline
Workers & Users & Req/s & Median & p95 \\
\hline
1 & 1 & 0.48 & 2.1\,s & 2.1\,s \\
1 & 4 & 1.21 & 3.3\,s & 3.5\,s \\
2 & 4 & 1.14 & 3.5\,s & 3.6\,s \\
4 & 4 & 1.04 & 3.8\,s & 4.2\,s \\
4 & 8 & 1.11 & 6.5\,s & 7.0\,s \\
\hline
\multicolumn{5}{l}{Worker kill: 0 failed probes, respawn 0.03\,s. Server crash: healthy in 0.9--1.0\,s.}\\
\multicolumn{5}{l}{Container restart: 3.8--3.9\,s. Fault-free probe: 60/60 over 5\,min.}\\
\end{tabular}
\end{table}
Throughput saturates near 1.2 clips/s on one node irrespective of worker count, because MediaPipe already parallelises
inside a request; scaling beyond one node was not tested. We make no availability claim: a 5-minute window cannot
support one.

% --------------------------------------------- V-C LIMITATIONS (replace)
\subsection{Limitations}
(1)~Eight athletes from a single session: the classification results are a
pilot and do not support claims of generalisation.
(2)~Tiers are cohort-relative tertiles of velocity; absolute m/s assumes a
1.70\,m stature (an error of about $\pm$5\% per 8.5\,cm of height difference).
(3)~2-D pose from a single side view at roughly 190\,px body height;
joint-angle accuracy was not validated against motion capture.
(4)~Two of nine watch files were not recorded during the filmed sprint.
(5)~Only the portion of the fly zone inside the camera's field of view
(0.9--1.7\,s) is analysed.
(6)~Latency was measured on one laptop CPU in two configurations, not on
phones.

% --------------------------------------------- REPRODUCIBILITY (new table)
\begin{table}[t]
\caption{Reproducibility settings}\label{tab:repro}
\centering\small
\begin{tabular}{ll}
\hline
Camera & Samsung Android phone %CONFIRM model
, 1080p HEVC, 60\,fps, tripod\\
Watch & Samsung Galaxy Watch %CONFIRM model
, Sensor Logger, 99.3--99.6\,Hz\\
Compute & Apple M5 (10 cores), 16\,GB RAM, macOS 27.0.1\\
Software & Python 3.11.17, MediaPipe 0.10.9, TF 2.18.1,\\
 & Keras 3.15.1, scikit-learn 1.7.2, NumPy 1.26.4,\\
 & SciPy 1.17.1, OpenCV 4.11, FFmpeg 9.0.2\\
Pose & complexity 1, det./track conf. 0.5\\
Seeds & 42 (MLP: 42--46)\\
Pipeline & DVC + MLflow, Docker image, \texttt{params.yaml}\\
 & MLflow 3.5.1, DVC 3.67.1, FastAPI 0.162.2, pandas 2.3.3\\
Code & \url{https://github.com/...} %CONFIRM\\
\hline
\end{tabular}
\end{table}

% --------------------------------------------- CONCLUSION (replace)
\section{Conclusion}
We presented a reproducible smartphone + smartwatch pipeline for sprint
biomechanics, with an objective, annotation-free velocity label, a
quantitatively validated video--IMU integration (stride-frequency MAE
0.07\,Hz, $r=0.95$; automatic detection of mismatched sensor files), and
measured per-stage latency (2.1\,s per clip; microsecond-scale classifier
inference in TFLite). A leakage-free pilot with 8 athletes showed that
three-tier skill classification is not yet reliable, and that earlier
reported accuracies were artefacts of label circularity and augmentation
leakage. A power analysis of KS-based drift monitoring sets the next
milestone: at least 100 labelled recordings before automated retraining
triggers become reliable. Because labelling is automatic, the pipeline can
reach that milestone at low cost through routine use.
